\documentclass[letterpaper]{article} % DO NOT CHANGE THIS
\usepackage[submission]{aaai2027}  % DO NOT CHANGE THIS
% The serif, sans-serif, and monospaced fonts are loaded automatically by
% aaai2027.sty (newtxtext, helvet, courier). DO NOT add \usepackage{times},
% \usepackage{helvet}, \usepackage{courier}, or any other font package.
\usepackage[hyphens]{url}  % DO NOT CHANGE THIS
\usepackage{graphicx} % DO NOT CHANGE THIS
\urlstyle{rm} % DO NOT CHANGE THIS
\def\UrlFont{\rm}  % DO NOT CHANGE THIS
\usepackage{natbib}  % DO NOT CHANGE THIS AND DO NOT ADD ANY OPTIONS TO IT
\usepackage{caption} % DO NOT CHANGE THIS AND DO NOT ADD ANY OPTIONS TO IT
\frenchspacing  % DO NOT CHANGE THIS
%
% These are recommended to typeset algorithms but not required.
\usepackage{algorithm}
\usepackage{algorithmic}
%
% Recommended for better-looking tables
\usepackage{booktabs}
%
% Keep the \pdfinfo as shown here.
\pdfinfo{
/TemplateVersion (2027.1)
}

\setcounter{secnumdepth}{1} % section numbers on (technical paper)

% ============================================================================
% TITLE
% 주의(익명성): "ARBOR"/"PySOAR" 라는 시스템 이름이 공개 저장소로 검색되면 double-blind
% 위배가 될 수 있음. 최종 제출 전 중립 명칭 사용 여부를 결정할 것.
% 대안 제목:
%  - Discovering Skills from Few Examples: An Impasse-Driven Symbolic Mechanism
%  - From Five Examples to a Skill: Program Synthesis and Anti-Unification over a
%    Symbolic Knowledge Graph
% ============================================================================
\title{ARBOR: Deriving Task-General Transformations from a\\
Handful of Examples by Impasse-Driven Search and Anti-Unification}

\author{
    Anonymous Submission
}
\affiliations{
    %
}

\begin{document}

\maketitle

\begin{abstract}
Humans acquire a new visual-reasoning skill from a handful of demonstrations,
whereas contemporary neural systems remain strikingly data-inefficient: even
frontier reasoning models score in the low single digits on novel few-example
abstraction tasks that ordinary people solve at once. We argue that the missing
ingredient is not more data or scale but a mechanism that \emph{derives}
reusable, verifiable structure from two to five examples. We present ARBOR, a
symbolic reasoning agent that treats each task as an analogy problem over a
five-layer knowledge graph and runs on a fidelity-first reimplementation of the
Soar decision cycle. Rather than executing a hand-written pipeline, ARBOR uses
general operators whose firing order is decided by \emph{impasse}: it descends
the knowledge-graph hierarchy only when it gets stuck, so different tasks yield
different reasoning traces. Within a training pair it synthesizes a program over
a deliberately minimal, frozen two-primitive space in which every argument is a
general expression \emph{discovered by search and verified on the training
pairs}; across pairs it aligns per-pair programs by a structure-mapping
correspondence and \emph{anti-unifies} them into a task-general schema with
typed variable slots. Consequently ``obvious'' notions---the largest object, the
bottom-right corner, coordinate formulas such as $(H-1,W-1)$---are never
hand-coded but emerge from a search whose tried-and-rejected hypotheses remain
auditable in the trace, and when the evidence is insufficient the agent declines
with an explicit impasse rather than silently fitting an answer. We describe the
mechanism and illustrate it on a controlled curriculum, reporting its scope
honestly; ARBOR is offered as a step toward data-efficient, interpretable skill
acquisition, not as a state-of-the-art benchmark solver.
\end{abstract}

% Uncomment to link code/data/extended version, but NOT for anonymous submission
% (a public repository link would de-anonymize the authors).
% \begin{links}
%     \link{Code}{...}
% \end{links}

% ============================================================================
\section{Introduction}
% [작성 완료 섹션] 흐름: (1) 인간 vs 기계 data efficiency + ARC → (2) 왜 스케일/ICL로
% 안 되는지(진단) → (3) 우리가 만드는 것(ARBOR 메커니즘) → (4) 기여 + 정직한 포지셔닝.
% ============================================================================

A person shown three or four examples of an unfamiliar grid transformation can
usually reproduce it on a new input immediately, inferring the underlying rule
rather than memorizing the pairs. This capacity for efficient skill acquisition
from minimal experience is what the Abstraction and Reasoning Corpus (ARC) was
designed to measure, deliberately controlling for prior knowledge so that raw
task accuracy cannot be ``bought'' with data or scale \cite{chollet2019measure}.
By this yardstick, the gap between human and machine learning remains wide.
Frontier reasoning systems that reach near-ceiling accuracy on the original
ARC-AGI-1 collapse to the low single digits on ARC-AGI-2, whose tasks require
composing several interacting rules from a few demonstrations, while typical
humans still solve them \cite{chollet2025arcagi2,arcprize2025}. Analyses of
large language models (LLMs) attribute this shortfall to a reliance on memorized
parametric knowledge rather than fluid, on-the-fly abstraction
\cite{wu2025fluid}, and even the newest interactive-reasoning benchmarks report
frontier models below one percent against a perfect human baseline
\cite{chollet2025arcagi2}. The bottleneck, in short, is not what these systems
know but how efficiently they can \emph{acquire and transfer} new structure.

Why does scale not close the gap? A growing body of work suggests that the
learning mechanism itself is the obstacle. In-context learning, often cited as
evidence of few-shot adaptation, behaves under distribution shift more like
selecting and tuning a function already contained in the pretraining hypothesis
space than like learning a genuinely new task \cite{wang2025icl}. Sequence
models attain near-perfect in-distribution accuracy yet fail on systematic
recombinations of familiar primitives \cite{lakebaroni2018scan,kim2020cogs}, and
minimal, meaning-preserving perturbations of a problem's surface form sharply
degrade even reasoning-tuned models, exposing template memorization rather than
rule extraction \cite{mirzadeh2025gsm}. Reinforcement-learning post-training,
meanwhile, has been shown to sharpen the base model's existing sampling
distribution rather than expand its reasoning boundary \cite{yue2025rl}. Symbolic
program induction offers the complementary strength---human-level generalization
from one or a few examples via compositional programs
\cite{lake2015human,ellis2021dreamcoder}---but existing systems pay for it with
stochastic wake--sleep search, large-scale sampling with replay, or per-task
gradient updates \cite{ellis2021dreamcoder,butt2024codeit,akyurek2024ttt}, and
recent library-learning methods lean on an LLM in the loop
\cite{grand2024lilo}. What is missing is a mechanism that extracts reusable,
train-verifiable structure from a handful of examples without a per-task
training loop.

We present \textbf{ARBOR}, a symbolic reasoning agent built to be exactly such a
mechanism. ARBOR treats every ARC task as an analogy problem and perceives it
only through ARCKG, a five-layer knowledge graph
(\textsc{task}$\rightarrow$\textsc{pair}$\rightarrow$\textsc{grid}$\rightarrow$%
\textsc{object}$\rightarrow$\textsc{pixel}) whose nodes expose typed properties.
It runs on a fidelity-first reimplementation of the Soar decision cycle
\cite{newell1990unified,laird2012soar}, differentially validated against a
reference Soar kernel. Crucially, ARBOR does not execute a scripted pipeline:
its general operators (observe, compare, hypothesize, synthesize, verify,
generalize, apply) are proposed by production rules and selected by the decision
cycle, and the agent descends the ARCKG hierarchy only when it reaches an
\emph{impasse}---an information gap it cannot resolve at the current level. As a
result the reasoning trace, including its descent depth and step count, differs
from task to task rather than following a fixed sequence. Within a single
training pair ARBOR synthesizes a program over a deliberately minimal, frozen
space of two primitives (a canvas constructor and a cell-coloring operator) in
which meaning lives not in choosing among many transformations but in
\emph{what expression each argument is}. Each argument is resolved as a general
expression discovered by a generate--apply--compare--reject search over
input-derived atoms and verified on the training pairs; thus a coordinate such
as the bottom-right corner $(H-1,W-1)$, or a selector such as ``the largest
object,'' is not assumed but \emph{emerges} from search, with the tried and
discarded hypotheses left visible in the trace. Across the training pairs ARBOR
aligns the per-pair programs by a structure-mapping correspondence and
\emph{anti-unifies} them---computing their least general generalization---into a
task-general schema whose differing positions become typed variable slots,
themselves resolved to input-derived expressions verified on training data only,
never on a test-label oracle. When the evidence underdetermines a slot, the
agent surfaces an explicit impasse and, in the retry setting, offers ranked
alternatives rather than silently committing to an unjustified answer.

\noindent\textbf{Contributions.} We make four contributions.
\begin{itemize}
\item \textbf{An impasse-driven mechanism, not a script.} We show how a faithful
Soar decision cycle over a symbolic knowledge graph produces a task-dependent
reasoning trajectory---descent depth and operator sequence emerge from
information gaps---so that a single set of general operators covers structurally
different tasks (Sections~3--4).
\item \textbf{Within-pair program synthesis with derived arguments.} Over a
frozen two-primitive space, ARBOR discovers each program argument as a general,
input-derived expression by search and train-verification, so notions such as
``largest,'' ``corner,'' and formulas like $(H-1,W-1)$ are found and checked
rather than hand-coded (Section~4).
\item \textbf{Cross-pair anti-unification for few-example generalization.} We
cast task-general skill acquisition as structure-mapping alignment followed by
anti-unification of per-pair programs into a schema with typed slots, resolved
from two to five examples without any per-task parameter update (Section~4).
\item \textbf{Honesty as a design and evaluation criterion.} The agent declines
via explicit impasse when information is insufficient, and its search is
auditable; we use per-task divergence of the reasoning trace, rather than
aggregate accuracy alone, as evidence that computation is genuine and not hidden
inside hand-written functions (Sections~5--6).
\end{itemize}

We stress the paper's scope. ARBOR is a \emph{mechanism} demonstrated on a
controlled curriculum of few-example tasks; it is not a state-of-the-art solver
for the full ARC-AGI corpus, and we report its coverage and failures plainly
(Section~6). Our aim is to make the components of efficient, transferable
abstraction explicit and inspectable, complementing the data-efficiency critique
of scaled neural learning with a concrete, interpretable alternative.

% ============================================================================
\section{Related Work}
% [작성 완료 섹션] 5개 소절: (A) LLM data efficiency 한계, (B) systematic/compositional
% generalization, (C) ARC/abstraction 벤치마크, (D) program synthesis & library
% learning + LLM-guided, (E) 유추·structure mapping·anti-unification + 인지아키텍처.
% 마지막에 ARBOR 차별화 문단.
% 주의: 아래 인용은 모두 실존 확인분 위주. arXiv 리포트는 bib에서 @misc.
% ============================================================================

\subsection{Data Efficiency and the Limits of Scaled Learning}
A recurring theme in recent analyses is that low test loss and strong
in-distribution accuracy do not by themselves yield the transferable,
sample-efficient behavior attributed to LLMs, and that a separate account beyond
statistical generalization is needed \cite{reizinger2024position}. Controlled
studies of in-context learning find that its apparent new-task learning holds
only in the absence of distribution shift, and that out-of-distribution it
resembles optimizing a function drawn from the pretraining hypothesis space
rather than acquiring a new one \cite{wang2025icl}. On mathematical reasoning,
perturbing only the surface form of otherwise identical problems degrades
even reasoning-oriented models and reveals a template-memorization failure mode
\cite{mirzadeh2025gsm}, and controllable-complexity puzzles show frontier
reasoning models collapsing beyond a modest difficulty threshold
\cite{shojaee2025illusion}. Reinforcement-learning post-training has likewise
been shown to reweight rather than extend a base model's reasoning support
\cite{yue2025rl}. These results motivate mechanisms that extract explicit,
verifiable structure rather than interpolate within a learned prior.

\subsection{Systematic and Compositional Generalization}
The difficulty of recombining known primitives in novel ways is longstanding.
The SCAN benchmark showed sequence-to-sequence networks failing on systematic
recombinations despite near-perfect training accuracy \cite{lakebaroni2018scan},
and COGS demonstrated a persistent structural-generalization gap in semantic
parsing even with ample data \cite{kim2020cogs}. Meta-learning has narrowed this
gap on controlled tasks \cite{lakebaroni2023}, but the challenge of reliably
producing systematic behavior from few examples remains central and directly
motivates a structure-mapping and anti-unification approach.

\subsection{Abstract Reasoning Benchmarks}
ARC operationalizes intelligence as skill-acquisition efficiency under strong
priors and few examples \cite{chollet2019measure}. Its successor ARC-AGI-2
raises the demand for multi-rule composition and symbol reuse, and reports
frontier reasoning systems far below human performance
\cite{chollet2025arcagi2,arcprize2025}. Analyses isolating fluid intelligence on
ARC trace LLM failure to weak skill composition, unfamiliarity with abstract
formats, and left-to-right decoding rather than to knowledge gaps
\cite{wu2025fluid}. We adopt this few-example, novel-task framing as the setting
our mechanism targets, while positioning our own evaluation on a controlled
curriculum rather than as a leaderboard claim.

\subsection{Program Synthesis and Library Learning}
Symbolic program induction achieves data-efficient generalization by
representing concepts as compositional programs, from Bayesian program learning
for one-shot concept acquisition \cite{lake2015human} to wake--sleep library
learning that discovers reusable abstractions from few-example tasks
\cite{ellis2021dreamcoder}. Subsequent work makes abstraction extraction fast
and corpus-guided \cite{bowers2023stitch}, adds LLM-guided synthesis with
symbolic compression and documentation \cite{grand2024lilo}, and refactors
LLM-generated programs into reusable libraries \cite{stengeleskin2024regal}.
A parallel line uses LLMs to propose programs or hypotheses for ARC-style tasks,
via natural-language hypothesis search \cite{wang2024hypothesis}, hindsight
replay for program search \cite{butt2024codeit}, or ensembling program induction
with direct transduction \cite{li2025induction}; test-time training closes much
of the gap through per-instance gradient updates on augmented examples
\cite{akyurek2024ttt}. ARBOR shares the goal of discovering reusable skills from
few examples but differs in its engine: a deterministic structure-mapping and
anti-unification core with no stochastic search loop, no LLM in the loop, and no
per-task parameter update.

\subsection{Analogy, Structure Mapping, and Cognitive Architectures}
Our mechanism is grounded in structure-mapping theory, in which an analogy maps
a system of relations by their syntactic structure rather than surface content,
preferring interconnected relational systems \cite{gentner1983structure}. Its
computational realization is anti-unification, the dual of unification that
computes a least general generalization capturing shared structure across
symbolic expressions \cite{cerna2023antiunification}. Melanie Mitchell's survey
frames abstraction and analogy-making as a central unsolved problem for AI,
noting that neither deep learning nor prior symbolic systems achieve human-like
abstraction \cite{mitchell2021abstraction}. Finally, ARBOR is built on the Soar
cognitive architecture and its impasse-driven, chunking-based problem solving
\cite{newell1990unified,laird2012soar}, which supplies the decision cycle that
turns our operators into a task-dependent reasoning trajectory rather than a
fixed script.

% ============================================================================
% 이하 §3-7 은 [스켈레톤]. 문단당 한 줄로 "무엇이 들어갈지"만 표시(값/성능 없음).
% 스토리텔링·구조 당위성 검토용. 실제 산문·수치는 나중에.
% 각 item = 계획된 한 문단. (These bullets are one-line-per-paragraph placeholders.)
% ============================================================================

\section{Background: The Soar Kernel and ARCKG}
\emph{[Skeleton --- one line per planned paragraph. Purpose: give the reader the
two substrates the mechanism needs, so Section~4 can be read as forced by them,
not chosen ad hoc. Ties back to Intro paragraph~3 and Related Work \S2.5.]}
\begin{itemize}
\item The reasoning substrate is a fidelity-first reimplementation of the Soar
decision cycle; we recover exactly the parts ARBOR depends on (preference-based
operator selection, impasses, truth maintenance, chunking).
\item We validate that kernel differentially against a reference Soar so that
the control layer is trustworthy rather than a bespoke loop.
\item The perceptual substrate is ARCKG, a five-layer typed-property knowledge
graph that is the agent's \emph{only} access to a task.
\item Access is lazy: the agent inspects a deeper layer only when it cannot
proceed at the current one, which is what makes descent meaningful rather than
routine.
\item The recursive \texttt{compare} operation supplies item-level and
relation-level (second- and higher-order) COMM/DIFF comparisons that are the sole
source of evidence downstream.
\end{itemize}

\section{The ARBOR Mechanism}
\emph{[Skeleton. Purpose: show each component is the minimal thing needed to turn
``few examples'' into a verified task-general transformation; organized by the
descriptive/procedural $\times$ pair/task quadrants. Delivers Intro
contributions~1--3; concrete alternative to Related Work \S2.4--2.5.]}

\subsection{Impasse-Driven Control}
\begin{itemize}
\item The agent's default is to attempt a task-general solution; failing that
raises an impasse that both diagnoses the missing information and triggers
descent --- so control is emergent, not scripted.
\item Because the trigger is an information gap, structurally different tasks
induce different operator sequences and descent depths from one rule set.
\end{itemize}

\subsection{Within-Pair Program Synthesis}
\begin{itemize}
\item Output space is deliberately frozen to two primitives (canvas plus
cell-coloring); expressive power is moved into \emph{what expression each
argument is}, not into a growing transformation catalogue.
\item Each argument is a general, input-derived expression found by a
generate--apply--compare--reject search and kept only if it reproduces every
training pair.
\item Hence selectors (``the largest object'') and coordinate formulas
(a bottom-right anchor) are \emph{discovered and verified}, never hand-coded, and
the tried-and-rejected candidates remain visible in the trace.
\item Surviving expressions are ranked by generality and simplicity so the agent
prefers the most transferable explanation consistent with the examples.
\end{itemize}

\subsection{Structure Mapping and Anti-Unification}
\begin{itemize}
\item Across pairs, per-pair programs are first aligned by a correspondence that
maximizes shared structure (structure mapping) --- the precondition for
generalizing.
\item Aligned programs are anti-unified position-by-position: agreeing parts stay
constant, differing parts become typed variable slots, yielding a task-general
schema.
\item Slots are resolved to input-derived expressions verified on the training
pairs only, with no appeal to a test-label oracle.
\item When the examples underdetermine a slot, the version space yields ranked
alternatives for the retry setting rather than a single unjustified commitment.
\end{itemize}

\subsection{Compression to Reusable Skills}
\begin{itemize}
\item When differing operation counts block anti-unification, pixel-level
programs are compressed into object-level ones, recovering a shared schema and
pointing toward reusable, accumulable skills.
\end{itemize}

\section{A Worked Example}
\emph{[Skeleton --- one line per planned paragraph; a single task carries the
whole pipeline so the reader sees the machinery actually run. Makes concrete the
Intro claim that ``obvious'' notions are derived, not assumed.]}
\begin{itemize}
\item Setup: one curriculum task and why it is a minimal but non-trivial instance
of the target capability.
\item Descent: the first impasse and the ARCKG path the agent takes to resolve
it, with the working-memory changes at each step.
\item Search: the coordinate/selector expressions tried and rejected before the
correct one survives --- the visible-search claim in action.
\item Generalization: the per-pair programs, their alignment, and the
anti-unified schema with its slots.
\item Application: instantiating the schema on the held-out input. (Figure~1:
trace/dashboard snapshot.)
\end{itemize}

\section{Demonstration and Honest Scope}
\emph{[Skeleton. Purpose: substantiate the mechanism honestly on a controlled
curriculum. No performance numbers here yet --- they are added later. Delivers
Intro contribution~4; fulfils the ``not a leaderboard'' stance of \S2.3.]}
\begin{itemize}
\item Curriculum: what the few-example task families are and what capability each
isolates.
\item Primary evidence is qualitative-structural: whether descent depth and
operator sequence \emph{diverge across tasks}, showing computation is real and
not hidden inside hand-written functions. (To be filled; no numbers yet.)
\item Interpretability: tried/rejected candidates and impasse-based declines are
inspectable in the trace. (Figure~2: dashboard.)
\item Scope stated plainly: task families covered, tasks declined, and the
boundary of current coverage --- with quantitative results deferred.
\item Contrast with scaled neural learning is drawn through the literature of
\S2, not through new LLM experiments.
\end{itemize}

\section{Limitations, Discussion, and Conclusion}
\emph{[Skeleton --- one line per planned paragraph.]}
\begin{itemize}
\item Limitation: coverage is bounded by the breadth of the expression space and
by task families outside the current operators.
\item Discussion: how compression toward reusable skills and richer relational
expressions widen coverage without abandoning the frozen-primitive discipline.
\item Outlook: extension to interactive settings (novel-environment adaptation
from few rules) as a natural next target for the same mechanism.
\item Conclusion: restate ARBOR as a concrete, interpretable, data-efficient
alternative to the learning the introduction diagnoses.
\end{itemize}

% ============================================================================
% References
% aaai2027.bst 사용(스타일은 aaai2027.sty 가 지정). 아래 arbor.bib 에 실존 확인된
% 항목만 수록. arXiv 리포트는 @misc(eprint/archivePrefix)로 표기(날조 venue 금지).
% ============================================================================
\bibliography{arbor}

\end{document}
